\section{The same Cartier cut with a fixed invertible coefficient}

\begin{definition}[The original fixed-coefficient power sequence]
\label{def:seshadri-fixed-coefficient-power-cut-sequence}
\lean{EnvelopeProperty.SeshadriAmplenessCriterionFixedCoefficientCutStability.fixedCoefficientPowerSectionSequence}
\leanok
\uses{def:seshadri-original-section-cut-complex,def:seshadri-same-section-original-power-complex}
Let $X$ be integral, let $L,E$ be original invertible sheaves, and let
$s:\mathcal O_X\to L^r$ be an original nonzero section. Put
$g:D\hookrightarrow X$ for its literal zero-subscheme inclusion.
For every $n\geq0$, define an actual sequence
\[
 E\otimes L^n\longrightarrow E\otimes L^{n+r}
 \longrightarrow g_*g^*(E\otimes L^{n+r}).
\]
Start with the existing section-cut sequence for $A=L^r$ and
$M=E\otimes L^n$. Transport its middle object by the canonical
isomorphism
\[
 c_n:A\otimes(E\otimes L^n)\simeq E\otimes L^{n+r}
\]
obtained from the inverse associator, the symmetry of $A,E$, the
associator, the symmetry of $A,L^n$, and the original power-addition
isomorphism. Transport the right object by $g_*g^*c_n$, leaving the
left object fixed. Thus both arrows come from the same original ideal
inclusion and quotient, with no supplied multiplication comparison.
\end{definition}
\begin{proof}
The existing section-cut construction has middle object
$A\otimes M$ and right object $g_*g^*(A\otimes M)$, with consecutive
composite zero. Conjugate the first arrow by the identity on $M$ and
$c_n$, and conjugate the second by $c_n^{-1}$ and $g_*g^*c_n$.
Their composite remains zero because $c_nc_n^{-1}$ is the identity.
These formulas specify the displayed actual sequence and its arrows.
\end{proof}

\begin{lemma}[Exactness with the same original coefficient]
\label{lem:seshadri-fixed-coefficient-power-cut-short-exact}
\lean{EnvelopeProperty.SeshadriAmplenessCriterionFixedCoefficientCutStability.fixedCoefficientPowerSectionSequence_shortExact}
\leanok
\uses{def:seshadri-fixed-coefficient-power-cut-sequence,lem:seshadri-actual-cartier-short-exact}
The sequence just defined is short exact for every $n\geq0$.
\end{lemma}
\begin{proof}
The original $E\otimes L^n$ is invertible, since both of its factors
are invertible. The existing Cartier-cut theorem therefore makes the
section-cut sequence for $A$ and $E\otimes L^n$ short exact.
The identity, $c_n$, and $g_*g^*c_n$ form an isomorphism of this short
complex with the specified fixed-coefficient sequence. Short exactness
is invariant under an isomorphism of short complexes, proving the claim.
\end{proof}

\begin{lemma}[Actual first-cohomology epimorphisms with fixed coefficient]
\label{lem:seshadri-fixed-coefficient-cut-eventual-hone-epi}
\lean{EnvelopeProperty.SeshadriAmplenessCriterionFixedCoefficientCutStability.exists_fixedCoefficient_singleCut_hOne_epi_of_ample}
\leanok
\uses{def:seshadri-fixed-coefficient-power-cut-sequence,lem:seshadri-fixed-coefficient-power-cut-short-exact,lem:vanishing-original-ample-serre,lem:seshadri-one-cut-cohomology-epi}
Let $X$ be integral and separated over a field $k$, equipped with an
original finite affine open cover. Keep the same original $L,E,r,s,D,g$.
If $g^*L$ is natively ample over $k$, then, for every sufficiently large
$n$, the actual first arrow of the fixed-coefficient sequence induces
an epimorphism
\[
 H^1(X,E\otimes L^n)\longrightarrow H^1(X,E\otimes L^{n+r}).
\]
\end{lemma}
\begin{proof}
Unpack native ampleness of $g^*L$ into its actual positive-power closed
embedding and whole power comparison. Apply the existing coherent Serre
theorem on the same $D$ with the fixed coherent coefficient $g^*E$.
It supplies one bound $N$ for vanishing of every positive cohomology
group of $g^*E\otimes(g^*L)^m$ whenever $m\geq N$.
The canonical tensor pullback and original power pullback isomorphisms
identify this sheaf with $g^*(E\otimes L^m)$, so the same groups vanish
for that literal original pullback. Since $n+r\geq n$, the same bound
applies when $n\geq N$ and $m=n+r$.

Apply $g^*c_n$ in reverse to obtain the vanishing on the actual right
term of the original section-cut sequence. The existing first-cohomology
epimorphism theorem for that sequence applies to $M=E\otimes L^n$.
Its first arrow followed by $c_n$ is exactly the specified
fixed-coefficient arrow. Functoriality of the original field-linear
cohomology and the isomorphism induced by $c_n$ preserve epimorphisms.
This proves the claim without assuming ambient ampleness or a target
vanishing statement.
\end{proof}

\begin{theorem}[The same fixed-coefficient first-cohomology maps stabilize]
\label{thm:seshadri-fixed-coefficient-cut-eventual-hone-iso}
\lean{EnvelopeProperty.SeshadriAmplenessCriterionFixedCoefficientCutStability.exists_fixedCoefficient_singleCut_hOne_isIso_of_ample}
\leanok
\uses{lem:seshadri-fixed-coefficient-cut-eventual-hone-epi,lem:seshadri-fixed-step-dimension-stabilization,thm:fc-scheme-coherent-euler-bounded,thm:fc-scheme-coherent-euler-zero}
In addition suppose $X$ is projective over $k$ and $r>0$. For every
sufficiently large $n$, the very same actual first-cohomology map is an
isomorphism. The bound may depend on the original fixed $E$.
\end{theorem}
\begin{proof}
Every $E\otimes L^n$ is invertible, hence a finite-type quasicoherent
module. Existing genuine projective coherent finiteness makes its
original field-linear first cohomology finite dimensional.
For $n$ beyond the preceding epimorphism bound, the target dimension is
therefore at most the source dimension. Apply the already proved
finite-residue stabilization lemma to these natural-number dimensions
with step $r$. It gives one larger bound at which the two dimensions
are equal. Each same actual map is then a surjection between
finite-dimensional vector spaces of equal dimension and hence an
injection. It is both monic and epic and is therefore an isomorphism.
The original finite-residue lemma is reused; no uniformity in $E$ or
incoming dimension-stability certificate is asserted.
\end{proof}

\begin{theorem}[The standard original restriction is onto with fixed coefficient]
\label{thm:seshadri-fixed-coefficient-cut-eventual-standard-restriction}
\lean{EnvelopeProperty.SeshadriAmplenessCriterionFixedCoefficientCutStability.exists_fixedCoefficient_singleCut_sections_surjective_of_ample}
\leanok
\uses{thm:seshadri-fixed-coefficient-cut-eventual-hone-iso,lem:seshadri-fixed-coefficient-power-cut-short-exact,lem:seshadri-actual-hone-mono-sections-onto,lem:seshadri-actual-cartier-quotient-standard-unit}
Under the preceding hypotheses, for every sufficiently large $n$ the
standard original adjunction-unit restriction
\[
 \Gamma(X,E\otimes L^{n+r})\longrightarrow
 \Gamma(D,g^*(E\otimes L^{n+r}))
\]
is surjective. In particular the fixed coefficient may be the actual
dual of a chosen original very ample line.
\end{theorem}
\begin{proof}
Use the actual short exact sequence of the first definition.
Its actual first-cohomology map is monic by the preceding theorem, so
the already proved natural degree-zero comparison and long exact
sequence make its original quotient map surjective on global sections.
The existing general Cartier quotient comparison, applied to
$A=L^r$ and $M=E\otimes L^n$, identifies its quotient with the standard
adjunction unit followed by pushforward of an internally constructed
automorphism of $g^*(A\otimes M)$. Conjugate this automorphism by
$g^*c_n$. Naturality of the same adjunction unit for $c_n^{-1}$ gives
the fixed-coefficient quotient as the standard unit on exactly
$E\otimes L^{n+r}$ followed by the conjugated automorphism.
The latter induces a bijection on global sections. Cancel this
bijection in the already proved quotient surjection. The resulting
surjection is precisely the stated standard restriction, with no
matching or surjectivity certificate supplied as input.
\end{proof}
